\documentclass[11pt,a4paper]{article}

\usepackage[T1]{fontenc}
\usepackage[utf8]{inputenc}
\usepackage{lmodern}
\usepackage{amssymb,amsmath,amsthm,mathtools}
\usepackage[margin=2.5cm]{geometry}
\usepackage{parskip}
\usepackage[hidelinks]{hyperref}
\usepackage{microtype}

% Theorem environments (arXiv / journal style)
\theoremstyle{plain}
\newtheorem{theorem}{Theorem}[section]
\newtheorem{proposition}[theorem]{Proposition}
\newtheorem{lemma}[theorem]{Lemma}
\newtheorem{corollary}[theorem]{Corollary}

\theoremstyle{definition}
\newtheorem{definition}[theorem]{Definition}
\newtheorem{example}[theorem]{Example}
\newtheorem{remark}[theorem]{Remark}

% Notation
\newcommand{\Ext}{\mathrm{Ext}}
\newcommand{\Shape}{\mathrm{Shape}}
\newcommand{\Pos}{\mathrm{Pos}}
\newcommand{\Type}{\mathbf{Type}}
\newcommand{\id}{\mathrm{id}}
\newcommand{\rootp}{\mathrm{o}}

\title{Covariant Functors of Directed Containers}
\author{Extracted from a machine-checked development}
\date{}

\begin{document}
\maketitle

\begin{abstract}
We introduce covariant functors between directed containers: maps that act
forward on both shapes and positions and preserve root, sub-shape, and shift.
Under the dictionary that a directed container presents a small category with
morphisms $p : s \to s \downarrow p$, such a functor is precisely a functor of
presented categories. We construct identities and composition, prove the
category laws, and record the transport lemmas needed for the composition law
on shifts. The proofs of all stated results have been machine-checked; we
record only the mathematical content.
\end{abstract}

\tableofcontents
\bigskip

%======================================================================
\section{Introduction}
%======================================================================

A directed container presents, for each shape~$s$, a category whose morphisms
out of~$s$ are positions of~$s$, with identity the root and composition given
by the shift. Container morphisms in the classical sense are
\emph{contravariant} on positions and correspond to cofunctors of these
presented categories. The present note isolates the covariant sibling: a map
that sends shapes forward and positions forward, preserving root, sub, and
shift.

The definition is given directly on directed containers, without passing
through an intermediate ``small category'' construction. This keeps the
identity and unit laws definitional on containers whose sub-shape operation is
strictly a projection, and matches the law-naming style of the corresponding
cofunctor development.

We work in a single universe of types. Proofs are indicated by the phrase
``by lean.''

\paragraph{What this note claims.}
Directed containers and their covariant functors form a category. A covariant
functor $C \Rightarrow D$ is a shape map and a forward position map satisfying
three laws F0--F2 (preservation of sub-shapes, roots, and shifts). Identities
and composition exist, composition is associative, and identities are strict
units.

%======================================================================
\section{Background: directed containers}
%======================================================================

\begin{definition}[Directed container]
A \textbf{directed container} consists of a container $(S,P)$ together with
\begin{align*}
  \rootp &\;:\; (s : S) \to P(s), \\
  \downarrow &\;:\; (s : S) \to P(s) \to S, \\
  \oplus &\;:\; (s : S) \to (p : P(s)) \to P(s \downarrow p) \to P(s),
\end{align*}
satisfying the standard axioms D1--D5.
\end{definition}

\begin{remark}[Presented category]
Informally, one may regard each shape~$s$ as an object and each position
$p \in P(s)$ as a morphism $s \to s \downarrow p$, with
$\id_s = \rootp(s)$ and composition $(p;q) = p \oplus q$. The axioms D1--D5
are the category laws for this presentation. We use this dictionary only as
motivation; the formal development never builds the category explicitly.
\end{remark}

%======================================================================
\section{Covariant functors}
%======================================================================

\begin{definition}[Covariant functor of directed containers]\label{def:dcont-functor}
Let $C$ and $D$ be directed containers. A \textbf{covariant functor}
$\varphi : C \Rightarrow D$ consists of
\begin{enumerate}
  \item a map on shapes $u : \Shape(C) \to \Shape(D)$, and
  \item for each shape $s \in \Shape(C)$, a map on positions
        \[
          g_s \;:\; \Pos(C)(s) \;\longrightarrow\; \Pos(D)\bigl(u(s)\bigr),
        \]
\end{enumerate}
such that the following three laws hold:
\begin{align}
  \tag{F0}
  u\bigl(s \downarrow_C p\bigr)
  &\;=\;
  u(s) \downarrow_D g_s(p),
  \\[4pt]
  \tag{F1}
  g_s\bigl(\rootp_C(s)\bigr)
  &\;=\;
  \rootp_D\bigl(u(s)\bigr),
  \\[4pt]
  \tag{F2}
  g_s\bigl(p \oplus_C q\bigr)
  &\;=\;
  g_s(p) \oplus_D
    \bigl(
      (\mathrm{F0}) \mathbin{\triangleright}
      g_{s \downarrow_C p}(q)
    \bigr).
\end{align}
Here the transport in F2 moves the inner image position
$g_{s \downarrow_C p}(q)$, which lives over $u(s \downarrow_C p)$, into the
fibre over $u(s) \downarrow_D g_s(p)$ along the equality F0.
\end{definition}

\begin{remark}
Contrast classical directed-container morphisms (cofunctors of the presented
categories), whose position component is \emph{contravariant}. Here both
components run forward: F0 says the codomain of a morphism is preserved, F1
preserves identities, and F2 preserves composition.
\end{remark}

%======================================================================
\section{Identity and transport lemmas}
%======================================================================

\begin{definition}[Identity functor]\label{def:id}
For a directed container~$C$, the \textbf{identity functor}
$\id_C : C \Rightarrow C$ has shape map and position map both the identity.
The laws F0--F2 hold by reflexivity; the transport in F2 is along a
reflexive equality and vanishes.
\end{definition}

\begin{lemma}[Position map commutes with transport]\label{lem:onpos-transport}
Let $\varphi : C \Rightarrow D$ be a covariant functor, let $e : a = b$ be an
equality of shapes of~$C$, and let $p$ be a position of~$a$. Then
\begin{equation}
  g_b\bigl(e \mathbin{\triangleright} p\bigr)
  \;=\;
  \bigl(u(e)\bigr) \mathbin{\triangleright} g_a(p),
\end{equation}
where $u(e)$ denotes the image of~$e$ under the shape map.
\end{lemma}

\begin{proof}
By lean.
\end{proof}

\begin{lemma}[Composite transport]\label{lem:transport-trans}
For a family $P$ and equalities $h_1 : x = y$, $h_2 : y = z$,
\begin{equation}
  h_2 \mathbin{\triangleright}
    \bigl(h_1 \mathbin{\triangleright} p\bigr)
  \;=\;
  (h_1 \mathbin{;}\, h_2) \mathbin{\triangleright} p.
\end{equation}
\end{lemma}

\begin{proof}
By lean.
\end{proof}

%======================================================================
\section{Composition}
%======================================================================

\begin{definition}[Composition of covariant functors]\label{def:comp}
Let $\varphi : C \Rightarrow D$ and $\psi : D \Rightarrow E$ be covariant
functors, with shape maps $u$, $u'$ and position maps $g$, $g'$ respectively.
Their \textbf{composite} $\psi \circ \varphi : C \Rightarrow E$ has
\begin{align}
  (\psi \circ \varphi)_S(s)
  &\;=\;
  u'\bigl(u(s)\bigr),
  \\[4pt]
  (\psi \circ \varphi)_p(s,p)
  &\;=\;
  g'_{u(s)}\bigl(g_s(p)\bigr).
\end{align}
The law F0 for the composite is the composite of the factors' F0 equalities;
F1 is two rewrites along the factors' F1; F2 is the factors' F2 followed by
Lemmas~\ref{lem:onpos-transport} and~\ref{lem:transport-trans}.
\end{definition}

\begin{proposition}[Left unit]\label{prop:id-comp}
For every covariant functor $\varphi : C \Rightarrow D$,
\begin{equation}
  \varphi \circ \id_C \;=\; \varphi.
\end{equation}
\end{proposition}

\begin{proof}
By lean.
\end{proof}

\begin{proposition}[Right unit]\label{prop:comp-id}
For every covariant functor $\varphi : C \Rightarrow D$,
\begin{equation}
  \id_D \circ \varphi \;=\; \varphi.
\end{equation}
\end{proposition}

\begin{proof}
By lean.
\end{proof}

\begin{proposition}[Associativity]\label{prop:comp-assoc}
For covariant functors
$\varphi : B \Rightarrow C$,
$\psi : C \Rightarrow D$, and
$\chi : D \Rightarrow E$,
\begin{equation}
  \chi \circ (\psi \circ \varphi)
  \;=\;
  (\chi \circ \psi) \circ \varphi.
\end{equation}
\end{proposition}

\begin{proof}
By lean.
\end{proof}

\begin{remark}
The data of functors compose strictly. The law fields are propositional, so
associativity and the unit laws hold by proof irrelevance once the underlying
maps agree.
\end{remark}

%======================================================================
\section{Discussion}
%======================================================================

\paragraph{What has been established.}
Covariant functors of directed containers form a category: they have
identities, an associative composition, and strict unit laws. The only
genuinely dependent obligation is F2, which requires transporting an inner
position along F0 before the target shift can consume it; composition of
functors inherits this obligation and discharges it via the two transport
lemmas.

\paragraph{Relation to cofunctors.}
The classical morphisms of directed containers (position-contravariant, shape
covariant) are the cofunctors of the presented categories. The structure
defined here is the functorial dual. Both live over the same objects; they are
not interchangeable.

\paragraph{Scope.}
The development does not construct the presented small category as an
independent object, nor does it prove an equivalence between covariant
functors of directed containers and functors of those categories. It also does
not treat natural transformations between covariant functors.

\paragraph{Machine checking.}
Every definition and theorem statement above corresponds to a declaration in a
machine-checked formalisation. The phrase ``by lean'' indicates that the proof
has been verified by a proof assistant.

%======================================================================
\section*{Acknowledgements}
%======================================================================

The mathematical content is a covariant counterpart to the cofunctor calculus
for directed containers of Ahman, Chapman, and Uustalu. The presentation here
is a prose rendering of a formal development intended for human review of the
theorem statements.

\begin{thebibliography}{9}

\bibitem{ahman2014}
D.~Ahman, J.~Chapman, and T.~Uustalu.
\newblock When is a container a comonad?
\newblock \emph{Logical Methods in Computer Science}, 10(3), 2014.

\bibitem{abbott2005}
M.~Abbott, T.~Altenkirch, and N.~Ghani.
\newblock Containers: constructing strictly positive types.
\newblock \emph{Theoretical Computer Science}, 342(1):3--27, 2005.

\end{thebibliography}

\end{document}
